\addto\captionsenglish{%
 \renewcommand\chaptername{Appendix}}
\appendixtitleon
\appendixtitletocon
\begin{appendices}
\chapter{Phụ lục: Chi tiết chiến lược giải bài toán MIOCP/MINLP}
\label{apd:A}

\section{Chiến lược giải bài toán MIOCP/MINLP}
\label{subsec:solver_strategy}

Bài toán MINLP \eqref{eq:joint_obj}--\eqref{eq:joint_binary} thuộc lớp NP-hard và
không có bộ giải thương mại hiệu suất cao tương đương Mosek/Gurobi cho lớp
MISOCP. Trong khuôn khổ này, chúng tôi sử dụng một chiến lược heuristic gồm ba bước: giải relaxation liên
tục tích hợp của MINLP sau rời rạc hóa, trích xuất đường đi rời rạc từ nghiệm
relaxation, rồi tinh chỉnh quỹ đạo bằng một NLP trên đường đi cố định. Các bài
toán NLP trong pipeline được xây dựng bằng CasADi và giải bằng IPOPT; các thao
tác trên đồ thị như tìm đường ngắn nhất hoặc liệt kê đường ứng viên được xử lý
riêng trên đồ thị vùng.

\paragraph*{Pha 1 — Relaxation liên tục tích hợp của MINLP.}
Tất cả biến của bài toán \eqref{eq:joint_obj}--\eqref{eq:joint_binary} được gom
vào một NLP duy nhất sau khi thả lỏng các biến nhị phân
$y_{uv},p_v\in\{0,1\}$ thành các biến liên tục trong đoạn $[0,1]$. Vector biến
quyết định gồm luồng cạnh $y_{uv}$, biến kích hoạt vùng $p_v$, trạng thái biên
$s_v^-,s_v^+$, tham số điều khiển $w_v$, thời lượng vùng $\Delta_v$, trạng thái
giao diện $z_{uv}$, và biến epigraph chi phí cục bộ $\rho_v$.

Cơ chế bật/tắt vùng được áp đặt bằng các ràng buộc Big-M dựa trên $p_v$. Cụ thể,
các biến trạng thái và thời lượng cục bộ được chặn bởi
\[
-M_s p_v \le s_v^\pm \le M_s p_v,\qquad
0\le \Delta_v \le M_\Delta p_v,\qquad
0\le \rho_v \le M_\rho p_v.
\]
Với tham số điều khiển, ta dùng dạng
\[
-M_w p_v \le w_v-\bar w_v \le M_w p_v,
\]
trong đó $\bar w_v$ là điều khiển danh nghĩa khi vùng không được kích hoạt. Khi
$p_v=0$, các biến động lực học của vùng bị tắt được đưa về giá trị danh nghĩa
và $\Delta_v=0$; do đó endpoint map thỏa
$F_v(0,\bar w_v,0)=0$, khiến ràng buộc defect
\eqref{eq:joint_defect} trở thành đẳng thức tầm thường mà không cần gắn thêm
Big-M trực tiếp vào defect. Trong cài đặt số, $F_v$ được xấp xỉ bằng tích phân
RK4 trên $N_{\text{int}}$ bước. Hàm mục tiêu được viết dưới dạng
$\min \sum_v \rho_v$ thông qua epigraph hóa chi phí cục bộ. Các hằng số Big-M
được chọn từ miền chặn vật lý của trạng thái, điều khiển, thời gian và chi phí;
giá trị quá lớn có thể làm bài toán NLP kém điều kiện số.

Ràng buộc hình học bật/tắt hỗ trợ hai chế độ:
\begin{itemize}
    \item \emph{Mesh-point hard constraint}: với polytope
    $Q_v=\{q:A_vq\le b_v\}$, ràng buộc
    $A_vq\le b_v+M_Q(1-p_v)$ được áp tại hai trạng thái biên $s_v^-,s_v^+$ và
    tại các điểm mesh nội tại của quỹ đạo RK4. Có thể thêm dung sai biên
    $\epsilon_b$ và khoảng dự phòng $\eta_{\text{safe}}$ để tránh nghiệm nằm quá
    sát biên vùng.
    \item \emph{Log-barrier penalty}: containment cứng vẫn được giữ tại hai
    trạng thái biên, còn an toàn nội tại được đưa vào hàm mục tiêu bằng số hạng
    \[
    -\mu \sum_v p_v\sum_{j}
    \log\!\bigl((b_v)_j+M_Q(1-p_v)-(A_vq)_j\bigr).
    \]
    Dịch slack $M_Q(1-p_v)$ được đặt bên trong logarit để vô hiệu hóa vùng không
    kích hoạt và tránh đánh giá logarit ngoài miền xác định.
\end{itemize}
Các ràng buộc ghép nối giao diện $s_u^+=z_{uv}=s_v^-$, điều kiện nguồn/đích, và
nếu cần cả ràng buộc liên tục điều khiển giữa hai vùng, được bật/tắt bằng Big-M
theo biến cạnh $y_{uv}$. Khi cạnh $(u,v)$ được kích hoạt, $z_{uv}$ đồng thời
phải nằm trong cả hai polytope $Q_u$ và $Q_v$.

\paragraph*{Khởi tạo (warm start).}
Để cải thiện khả năng hội tụ của IPOPT, nghiệm khởi tạo được xây dựng từ một
đường đi hình học $\pi_0$ trên đồ thị vùng. Cụ thể, $\pi_0$ là đường ngắn nhất
từ $\sigma$ đến $\tau$ theo trọng số khoảng cách giữa tâm các vùng. Với mỗi vùng
$v\in\pi_0$, các trạng thái $s_v^-$ và $s_v^+$ được đặt tại các điểm neo trên
giao của hai vùng kề nhau; nếu giao bị thoái hóa, dùng trung điểm giữa hai tâm
vùng làm xấp xỉ. Góc hướng $\theta$ được suy từ vector nối điểm vào và điểm ra,
$\Delta_v$ được ước lượng từ chiều dài đoạn chia cho vận tốc tối đa, còn $w_v$
được đặt bằng điều khiển danh nghĩa. Các biến luồng được khởi tạo bằng
$y_{uv}=1,p_v=1$ trên $\pi_0$ và bằng $0$ trên các cạnh, vùng còn lại.

\paragraph*{Continuation trên tham số rào.}
Khi sử dụng log-barrier, chúng tôi áp dụng kỹ thuật continuation trên tham số
$\mu$~\cite{BoydVandenberghe2004}. Cụ thể, giải một dãy NLP với $\mu$ giảm dần
từ $\mu_0=10^{-2}$ về $\mu_{\min}=10^{-5}$, dùng nghiệm ở bước trước làm khởi
tạo cho bước tiếp theo. Số vòng continuation và hệ số giảm $\mu$ là các tham số
cấu hình. Với chế độ mesh-point hard constraint, bước continuation này được bỏ
qua.

\paragraph*{Pha 2 — Trích xuất đường đi rời rạc.}
Sau khi NLP relaxation hội tụ, nghiệm $y^\star\in[0,1]^{|E|}$ được chuyển thành
một đường đi rời rạc $\pi^\star$. Trước tiên, ta duyệt tham lam từ $\sigma$:
tại mỗi đỉnh, chọn cạnh ra có $y_{uv}^\star$ lớn nhất nếu giá trị này vượt
ngưỡng $0.5$; nếu không có cạnh nào đạt ngưỡng, hạ ngưỡng xuống $10^{-3}$ để
tránh dừng sớm. Nếu thủ tục tham lam không đến được $\tau$, hoặc độ lệch nguyên
\[
\delta_{\mathrm{int}}
=\max_{e\in E}\min\{y_e^\star,\,1-y_e^\star\}
\]
lớn hơn $10^{-3}$, ta thay bằng đường đi ngắn nhất trên đồ thị con
$\{e:y_e^\star>10^{-6}\}$ với trọng số $-\log y_e^\star$. Cách làm này ưu tiên
các cạnh có giá trị luồng lớn trong nghiệm relaxation, nhưng vẫn trả về một
đường đi rời rạc hợp lệ.

\paragraph*{Pha 3 — Polish quỹ đạo trên đường đi cố định.}
Nghiệm relaxation có thể còn phân số, nên các trạng thái nối $s_u^+$ và $s_v^-$
trên hai vùng liên tiếp của $\pi^\star$ có thể chưa khớp hoàn toàn. Ta đo sai
lệch này bằng \emph{connection gap}. Nếu
$\max\{\delta_{\mathrm{int}},\delta_{\mathrm{conn}}\}$ vượt ngưỡng cấu hình,
hoặc nếu muốn loại bỏ hoàn toàn ảnh hưởng của Big-M, ta cố định đường đi
$\pi^\star$ và giải lại một NLP \emph{fixed-path}. Bài toán này chỉ giữ các biến
liên tục trên các vùng thuộc $\pi^\star$; các ràng buộc ghép nối giao diện
trở thành đẳng thức cứng, kèm điều kiện biên start/goal, ràng buộc defect
multiple shooting, và ràng buộc containment theo cùng chế độ ở Pha 1. Đây là
dạng direct multiple shooting tiêu chuẩn trên một chuỗi vùng đã biết
~\cite{BockPlitt1984, DiehlEtAl2006, Betts2010}. Nghiệm relaxation được dùng
làm khởi tạo cho NLP fixed-path; nếu dùng log-barrier, continuation trên $\mu$
được tiếp tục áp dụng trong bước polish.

% \paragraph*{Fallback bằng các đường đi ứng viên.}
% Nếu polish trên $\pi^\star$ thất bại, ví dụ IPOPT không hội tụ hoặc residual của
% ràng buộc vượt ngưỡng cho phép, thuật toán xét thêm tối đa $K$ đường đi đơn ứng
% viên từ $\sigma$ đến $\tau$. Các đường này được liệt kê theo thứ tự tăng dần của
% chi phí hình học bằng thuật toán Yen~\cite{Yen1971}; trong cài đặt, thủ tục này
% được thực hiện bằng \texttt{networkx.shortest\_simple\_paths}. Với từng đường
% ứng viên, ta giải NLP fixed-path như ở Pha 3 và chọn nghiệm khả thi có chi phí
% thấp nhất trong số các nghiệm hội tụ. Tham số $K$ kiểm soát trực tiếp đánh đổi
% giữa độ tin cậy của fallback và thời gian tính toán.

% \paragraph*{Lý do không dùng bộ giải MINLP đơn khối.}
% Các bộ giải MINLP đa năng như Bonmin hay Juniper~\cite{LubinEtAl2018} thường
% thực hiện branch-and-bound trên biến nhị phân và giải một NLP tại mỗi nút của
% cây tìm kiếm~\cite{Wolsey1999}. Cách tiếp cận này đúng về mặt mô hình hóa, nhưng
% không phù hợp với mục tiêu tính toán của bài toán hiện tại. Thứ nhất, việc giải
% lặp lại một NLP lớn tại nhiều nút có chi phí cao hơn đáng kể so với giải một
% relaxation tích hợp rồi chỉ polish trên một số ít đường đi cố định. Thứ hai,
% phần rời rạc của bài toán có cấu trúc luồng đơn vị từ $\sigma$ đến $\tau$; khi
% relaxation đã cung cấp thông tin tốt về các cạnh có khả năng được chọn, việc
% trích xuất và kiểm tra một số đường ứng viên thường rẻ hơn so với khám phá toàn
% bộ cây branch-and-bound.

% Chiến lược trên không đảm bảo tối ưu toàn cục. IPOPT chỉ tìm nghiệm dừng địa
% phương cho các NLP phi lồi, và đường đi tối ưu toàn cục của MINLP có thể không
% được chọn bởi thủ tục trích xuất hoặc không nằm trong tập $K$ đường ứng viên.
% Do đó, nghiệm cuối cùng được hiểu là nghiệm khả thi hoặc tối ưu địa phương trong
% tập đường đi đã xét. Đây là đánh đổi có chủ đích: giảm bảo đảm tối ưu toàn cục
% để đổi lấy thời gian giải khả thi trên các đồ thị có hàng chục đến hàng trăm
% vùng. Chất lượng nghiệm, ảnh hưởng của $K$, các ngưỡng
% $\delta_{\mathrm{int}}$, $\delta_{\mathrm{conn}}$, chế độ containment, và thời
% gian tính toán sẽ được đánh giá trong chương thực nghiệm.


\end{appendices}